% !TEX root = ../Main.tex
\chapter{蝶泳}

\section{身体姿势}

\state{俯卧与禁止动作}{
    \begin{itemize}
        \item 从出发和每次转身后的\textbf{第一次手臂动作}开始，身体应保持\textbf{俯卧姿势}；
        \item 允许\textbf{水下侧打腿}；
        \item \textbf{任何时候都不允许呈仰卧姿势}。
    \end{itemize}
}

\section{手臂动作}

\state{双臂同时的摆臂与划水}{
    \begin{itemize}
        \item 两臂必须\textbf{在水面上同时向前摆动}；
        \item 并\textbf{同时在水下向后划水}。
    \end{itemize}
}

\section{腿部动作}

\state{上下打水的同时性}{
    \begin{itemize}
        \item 所有腿部的\textbf{上下打水动作必须同时进行}；
        \item 两腿或两脚\textbf{可不在同一水平面上}；
        \item 但\textbf{不允许有交替动作}；
        \item \textbf{不允许蹬蛙泳腿}。
    \end{itemize}
}

\section{转身与终点}

\state{双手同时触壁}{
    在每次\textbf{转身}和\textbf{到达终点}时，\textbf{两手应在水面、水上或水下同时触壁}。
}

\section{出发与 15 米规则}

\state{出发 / 转身后的水下动作}{
    \begin{itemize}
        \item 出发和每次转身后，允许运动员身体\textbf{完全没入水中}；
        \item 可做\textbf{一次或多次打腿动作}和\textbf{一次划水动作}；
        \item 划水动作结束时，必须使身体\textbf{升出水面}；
        \item 在\textbf{15 米前（含 15 米）}，运动员的\textbf{头部必须露出水面}；
        \item 运动员必须使身体\textbf{保持在水面上}，直至\textbf{下次转身或到达终点}。
    \end{itemize}
}

\rmkb{
    \textbf{易考点}：
    \begin{itemize}
        \item 蝶泳始终\textbf{俯卧}（绝不仰卧）；允许水下\textbf{侧打腿}；
        \item 两臂\textbf{同时}在水面前摆、\textbf{同时}在水下后划；
        \item 两腿\textbf{上下打水同时进行}，可不在同一水平面但\textbf{不得交替}、\textbf{不可蹬蛙泳腿}；
        \item \textbf{双手同时触壁}（转身与终点）；
        \item 15 米规则：水下允许多次打腿 $+$ 一次划水，15 米前（含）头必须露出。
    \end{itemize}
}
